\newpage

\section[Cap\'{i}tulo 3: Gamificaci\'{o}n y Juegos Serios]{\centering Cap\'{i}tulo 3: Gamificaci\'{o}n y Juegos Serios}

La gamificaci\'{o}n y los juegos serios permiten explicar c\'{o}mo una experiencia interactiva puede apoyar procesos formativos sin convertirse en una clase tradicional dentro de una pantalla. En este proyecto, estos conceptos ayudan a dise\~{n}ar escenarios donde el jugador aprende al interpretar informaci\'{o}n, tomar decisiones y revisar consecuencias (Deterding et al., 2011).

La referencia a videojuegos como \textit{Papers, Please} y \textit{Return of the Obra Dinn} se entiende como una inspiraci\'{o}n de mec\'{a}nicas, no como reproducci\'{o}n de su contenido. Lo relevante para el proyecto es la atenci\'{o}n a detalles, la detecci\'{o}n de incongruencias y la reconstrucci\'{o}n de hechos a partir de pistas presentes en pantalla (Michael \& Chen, 2006).

\subsection{Gamificaci\'{o}n}

La gamificaci\'{o}n consiste en incorporar elementos propios del juego en contextos que no son juegos completos, como puntos, progresi\'{o}n, retos, niveles o recompensas. Su valor no est\'{a} solamente en hacer una actividad m\'{a}s entretenida, sino en orientar la participaci\'{o}n del usuario hacia objetivos concretos (Deterding et al., 2011).

En el proyecto, la gamificaci\'{o}n se refleja en niveles, casos progresivos, reportes de desempe\~{n}o y consecuencias narrativas. Estos elementos buscan mantener la motivaci\'{o}n del estudiante mientras practica habilidades de observaci\'{o}n y respuesta ante incidentes simulados (Kapp, 2012).

\subsection{Juegos serios}

Los juegos serios son experiencias dise\~{n}adas con un prop\'{o}sito formativo, informativo o de entrenamiento, aunque mantengan caracter\'{i}sticas propias del entretenimiento. En ciberseguridad, este enfoque permite representar errores y riesgos de manera segura, sin exponer sistemas reales (Michael \& Chen, 2006).

\begin{center}
\begin{tabular}{|p{0.30\textwidth}|p{0.55\textwidth}|}
\hline
\textbf{Elemento} & \textbf{Uso dentro del proyecto} \\
\hline
Escenario simulado & Ambiente donde el jugador analiza correos, alertas o incidentes sin afectar sistemas reales (Ng \& Hasan, 2024). \\
\hline
Reglas de decisi\'{o}n & Acciones disponibles seg\'{u}n el tipo de amenaza y el nivel de riesgo observado (Kapp, 2012). \\
\hline
Consecuencias & Reportes posteriores que muestran efectos de las decisiones tomadas durante el nivel (Hattie \& Timperley, 2007). \\
\hline
Progresi\'{o}n & Aumento gradual de dificultad mediante casos con m\'{a}s se\~{n}ales, ruido e informaci\'{o}n parcial (Biggs \& Tang, 2011). \\
\hline
\end{tabular}
\end{center}

\subsection{Relaci\'{o}n con la ense\~{n}anza de ciberseguridad}

En ciberseguridad, los juegos serios resultan pertinentes porque permiten representar ataques, errores humanos y procesos de respuesta sin comprometer infraestructura real. Adem\'{a}s, facilitan que el estudiante repita escenarios y mejore su criterio mediante pr\'{a}ctica guiada por consecuencias (Ng \& Hasan, 2024).

La propuesta adopta estas ideas para construir un prototipo centrado en observaci\'{o}n, comparaci\'{o}n y decisi\'{o}n. El objetivo no es entregar respuestas obvias, sino crear situaciones donde el estudiante deba identificar pistas y comprender por qu\'{e} una acci\'{o}n resulta m\'{a}s adecuada que otra (Irvine et al., 2005).

\subsection{S\'{i}ntesis del cap\'{i}tulo}

La gamificaci\'{o}n aporta estructura de avance y motivaci\'{o}n, mientras que los juegos serios aportan el enfoque formativo. En conjunto, ambos conceptos permiten plantear un videojuego que ense\~{n}a de forma pr\'{a}ctica, usando escenarios y consecuencias en lugar de explicaciones directas.
